\documentclass[10pt]{article}
\usepackage[margin=0.85in]{geometry}
\usepackage{amsmath,amssymb}
\begin{document}
\begin{center}
  {\Large\bfseries A quadratic vertex bound for exact DP-coloring gaps\par}
  \smallskip
  Disproof of Conjecture 00000003478
\end{center}

For every integer \(g\ge 0\), there is a bipartite graph on at most
\(4(g+1)^2\) vertices with \(\chi_{\mathrm{DP}}-\chi=g\). This disproves
the exponential growth assertion about the minimum vertex count.

\paragraph{Cyclic correspondence covers.}
In a \(k\)-fold DP cover, each vertex has \(k\) available colors, and each
edge forbids a matching of pairs of colors at its endpoints. A coloring
selects one color per vertex and avoids the forbidden pair on every edge.
The DP chromatic number is the smallest \(k\) for which every such cover
is colorable. This is the usual correspondence-coloring definition
\cite{bernshteyn}. If every \(k\)-fold cover is colorable, the same holds
for larger cover sizes by restricting every color set to \(k\) colors.

Fix \(k\ge2\) and put \(m=2k^2\). On the complete bipartite graph
\(K_{m,m}\), use the color set \(\mathbb Z/k\mathbb Z\) at every vertex.
For each edge \(uv\), directed from the left side to the right side,
choose a shift \(s_{uv}\in\mathbb Z/k\mathbb Z\). Forbid the pairs
\((a,a+s_{uv})\). Each shift is a perfect matching and the shifts can
be chosen independently, yielding \(k^{m^2}\) covers. For a fixed
assignment of colors to the \(2m\) vertices, exactly one shift violates
the assignment on each edge, so exactly \((k-1)^{m^2}\) covers admit that
assignment. Thus at most \(k^{2m}(k-1)^{m^2}\) covers admit some assignment.

The first two terms of the binomial expansion give
\[
 k^k=((k-1)+1)^k\ge (k-1)^k+k(k-1)^{k-1}>2(k-1)^k.
\]
Since \(k\le 2^k\) and \(m^2=k(2km)\), it follows that
\[
 k^{2m}(k-1)^{m^2}
 \le 2^{2km}(k-1)^{m^2}
 =\bigl(2(k-1)^k\bigr)^{2km}
 <(k^k)^{2km}=k^{m^2}.
\]
Some \(k\)-fold cover therefore has no coloring, and
\(\chi_{\mathrm{DP}}(K_{m,m})\ge k+1\).

\paragraph{Attaining any exact gap.}
For a vertex \(v\) of a graph \(G\), one has
\[
 \chi_{\mathrm{DP}}(G-v)\le \chi_{\mathrm{DP}}(G)
 \le \chi_{\mathrm{DP}}(G-v)+1.
\]
The first inequality follows by extending a cover of \(G-v\) to \(G\)
with arbitrary colors at \(v\) and empty matchings at edges incident
with \(v\). For the second, take \(q=\chi_{\mathrm{DP}}(G-v)\) and an
arbitrary \((q+1)\)-fold cover of \(G\). Choose a color at \(v\), then
delete the color paired with it at each neighbor if such a pair exists.
Every remaining vertex has at least \(q\) available colors. Restrict to
\(q\) colors at each and color \(G-v\). The prescribed color at \(v\)
extends this to a coloring of the original cover.

Delete vertices of \(K_{m,m}\) one at a time, leaving at least one on each
side, until only \(K_{1,1}\) remains. Every intermediate graph has
ordinary chromatic number two. At the end the DP chromatic number is
two as well. A forbidden pair can block a one-fold cover of the last edge.
In any two-fold cover, a color at one endpoint forbids at most one of the
two colors at the other. The DP chromatic numbers can only stay constant or decrease
by one at each step. For \(g\ge1\), take \(k=g+1\), so the sequence starts
at DP chromatic number at least \(g+2\) and ends at two. It therefore
contains a value exactly \(g+2\). This yields a graph of gap exactly \(g\)
on at most \(2m=4(g+1)^2\) vertices. For \(g=0\), take \(K_2\) itself.
The resulting quadratic upper bound is incompatible with a lower bound
exponential in \(g\).

\paragraph{Lean verification.}
The Lean model encodes a cover as a word of length \(m^2\) and an assignment
as a word of length \(2m\). It certifies the row-major edge indices and
checks the matching property and uniqueness of the forbidden shift
modulo \(k\), counts the full cover list without repetition, and counts the
concatenated lists of covers admitted by each assignment. The theorem
\texttt{exists\_uncolorable\_cyclic\_cover} uses the strict counting inequality
to produce a cover outside this concatenation. The theorem
\texttt{exact\_gap\_of\_deletion} verifies the discrete hitting and vertex-bound
argument from the graph deletion facts proved above.

\begin{thebibliography}{1}
\bibitem{bernshteyn} A. Bernshteyn, \emph{The asymptotic behavior of the
correspondence chromatic number}, Discrete Math. \textbf{339} (2016),
2680--2692.
\end{thebibliography}
\end{document}
